\documentclass[10pt]{article}
\usepackage[margin=0.88in]{geometry}
\usepackage{amsmath,amssymb,booktabs,array}
\usepackage[hidelinks]{hyperref}
\usepackage[T1]{fontenc}
\usepackage{lmodern}
\setlength{\parindent}{0pt}
\setlength{\parskip}{6pt}
\title{Supporting Information\\[2pt]\large When Does Posterior Integration Change Option Values? Evidence from WTI Average Price Options}
% Deliberately blinded: submit the identifying title page separately.
\author{}
\date{}
\begin{document}
\maketitle
\noindent\textbf{Scope.}
This supplement reports ancillary design choices, error benchmarks,
simulation diagnostics, liquidity checks, and numerical sensitivity for the
accompanying research article. Synthetic model checks and observed
market-price-proxy comparisons are separate sources of evidence.
Dollar prices are per barrel unless stated otherwise.

\section*{S1. Historical-volatility input benchmarks}
The main conditional comparison contains 2,381 APO contract-date observations
from 15 valuation dates. Model repricing uses the valuation-date futures
curve and known fixings; only the option-implied volatility input is
estimated from strictly earlier option dates. Results should not be
interpreted as a fully prospective price forecast.

\begin{table}[h]
\centering
\caption{Table S1. Conditional pricing errors (USD/bbl).}
\small
\begin{tabular}{lrrrr}
\toprule
Volatility input & $N$ & Dates & MAE & RMSE\\
\midrule
Long-window historical PI & 2,381 & 15 & 2.6187 & 3.4090\\
Rolling 252 returns & 2,381 & 15 & 5.2183 & 6.3933\\
Rolling 126 returns & 2,381 & 15 & 8.3419 & 9.7846\\
Rolling 63 returns & 2,381 & 15 & 4.5201 & 5.6465\\
EWMA (63-day half-life) & 2,381 & 15 & 5.6616 & 6.8824\\
Horizon-matched GARCH(1,1) & 2,381 & 15 & 2.6730 & 3.4441\\
Earlier-date expanding APO smile & 2,381 & 15 & 0.1088 & 0.1725\\
Previous-day APO smile & 2,381 & 15 & 0.1075 & 0.1594\\
\bottomrule
\end{tabular}
\end{table}

GARCH parameters are fitted using prior returns; the valuation-date
return updates conditional variance without re-estimating the parameters.
Expected daily variances are forecast over the unresolved fixing
horizon, then mapped to the single constant-volatility parameter of the
maintained one-factor price representation. The comparison thus
uses some same-day state information.

\section*{S2. Numerical pricing checks and posterior calibration}
The within-APO PI--PM differences are small. For representative October
calls at strike 98, independent pseudo-random Monte Carlo and scrambled
Sobol simulations of the payoff agree with Curran conditioning about the
\emph{PI--PM gap}, not necessarily the absolute theoretical option price.
Pricing at posterior root-mean-square volatility
$\sigma_{\mathrm{RMS}}=\sqrt{\mathbb E(\sigma^2\mid\mathcal D)}$
is an additional plug-in sensitivity, not posterior integration.

\begin{table}[h]
\centering
\caption{Table S2A. Representative PI--PM gaps (USD/bbl).}
\small
\begin{tabular}{lrrr}
\toprule
Date / call & Pseudo-MC & Curran & Scrambled Sobol\\
\midrule
2026-08-27, Oct., $K=98$ & 0.0013683 & 0.0013690 & 0.0013717\\
2026-08-28, Oct., $K=98$ & 0.0013394 & 0.0013402 & 0.0013393\\
\bottomrule
\end{tabular}
\end{table}

\begin{table}[h]
\centering
\caption{Table S2B. Prior-predictive Gaussian posterior simulation calibration.}
\small
\begin{tabular}{rrrr}
\toprule
History size $n$ & Mean rank & KS $p$-value & 95\% coverage\\
\midrule
21 & 0.5000 & 0.861 & 0.9500\\
63 & 0.4993 & 0.348 & 0.9505\\
252 & 0.4978 & 0.210 & 0.9483\\
1,260 & 0.5027 & 0.106 & 0.9505\\
\bottomrule
\end{tabular}
\end{table}

The simulation-based calibration study contains 200,000 prior-predictive
datasets, 50,000 for each history size. It is a check on a particular
posterior implementation under its assumed model; it does not validate
real WTI stochastic dynamics.

\section*{S3. Liquidity and valuation-date dependence}
Of the 2,381 date-ordered contract observations, 2,319 (97.4\%)
report zero daily volume. Comparisons resample whole valuation dates,
preserving within-date cross-option dependence but not removing serial
dependence across adjacent dates. The positive-volume subset is small.

\begin{table}[h]
\centering
\caption{Table S3. Positive-volume conditional comparisons.}
\small
\begin{tabular}{lrrrr}
\toprule
Volatility input & $N$ & Dates & MAE & RMSE\\
\midrule
Historical PI & 62 & 13 & 1.6340 & 2.3391\\
Earlier-date APO expanding smile & 62 & 13 & 0.1003 & 0.1358\\
Previous-day APO smile & 62 & 13 & 0.1084 & 0.1333\\
\bottomrule
\end{tabular}
\end{table}

The same aggregate ordering occurs at positive open-interest
thresholds of 10, 100, and 500 contracts. At open interest of at
least 100, the APO implied-volatility MAEs are approximately
0.106--0.107, versus 1.990 for the historical benchmark. Such
subsamples do not establish executable tradability or unbiased
transaction-price discovery.

\section*{S4. Cross-instrument WTI volatility transfer}
Monthly American WTI (LO) futures-option implied volatilities are
converted to a scalar APO volatility using the fixing-count
weighted RMS mapping. An alternative mapping matches the implied
variance of the unresolved arithmetic-average component under the
maintained common-factor representation.

\begin{table}[h]
\centering
\caption{Table S4A. LO cross-instrument mappings on 253 common-support states.}
\small
\begin{tabular}{lrr}
\toprule
Specification & MAE & RMSE\\
\midrule
Previous-day LO, weighted RMS & 0.1559 & 0.2599\\
Previous-day LO, variance-matched & 0.1562 & 0.2612\\
Expanding LO, weighted RMS & 0.1604 & 0.2858\\
Expanding LO, variance-matched & 0.1611 & 0.2869\\
\bottomrule
\end{tabular}
\end{table}

\begin{table}[h]
\centering
\caption{Table S4B. Absolute LO implied-volatility changes under tree refinement.}
\small
\begin{tabular}{lrrrr}
\toprule
Tree steps & Mean & Median & 95th pct. & Maximum\\
\midrule
80 to 160 & 0.000820 & --- & --- & ---\\
160 to 320 & 0.000402 & 0.000284 & 0.001089 & 0.001810\\
\bottomrule
\end{tabular}
\end{table}

The refinement audit uses 156 stratified LO contract-dates.
The 160-step baseline is practically stable for this comparison;
the audit does not prove the infinite-step tree limit. A
strictly prior smile-shape audit examined 144 diagnostics over 12
dates. Fitted volatility was finite and strictly positive within
observed prior moneyness support, with at most one direction change.

\section*{S5. Bounded synthetic model-risk and grid sensitivity}
The pilot-informed independent-seed study contains 160 distinct
historical Gaussian return draws, 40 each at $(\sigma,n)$ values
$(0.20,21)$, $(0.20,252)$, $(0.50,21)$, and $(0.50,252)$.
Three strikes (95, 100, 105) share each return history. For
each state, the same 95\% posterior price interval is assessed
against three independently simulated known-volatility payoff
references, with hypothetical cross-contract correlation weights
$\rho=1,0.75,0.50$. No correlation stress changes the posterior
pricing rule. Therefore 1,440 option-scenario rows are not
statistically independent.

An independent scrambled Sobol payoff calculation with exact
geometric-average control supplies the fixed-parameter payoff
references. Twelve independently scrambled 65,536-pair blocks
at $\rho=1$ and ten 16,384-pair blocks for each stress correlation
provide Monte Carlo reference uncertainty estimates. The largest
one-factor reference 95\% halfwidth is \$0.000154/bbl, and the
largest observed absolute Curran approximation discrepancy in six
one-factor controls is \$0.000922/bbl.

\begin{table}[h]
\centering
\caption{Table S5. Full 1,440-row posterior interval grid check:
coverage counts do not change across all tested grids.}
\small
\begin{tabular}{rrrrr}
\toprule
Volatility grid & Coverage $\rho=1$ & $\rho=0.75$ & $\rho=0.50$ &
Flipped decisions\\
\midrule
513 & 95.00\% & 65.21\% & 31.88\% & 0\\
1,025 & 95.00\% & 65.21\% & 31.88\% & 0\\
2,049 & 95.00\% & 65.21\% & 31.88\% & 0\\
4,097 & 95.00\% & 65.21\% & 31.88\% & 0\\
\bottomrule
\end{tabular}
\end{table}

The 513-point results exactly reproduce all 1,440 stored rows
(including PI, PM and both posterior price-interval endpoints).
Relative to 4,097 points, the maximum endpoint displacement was
\$0.005227/bbl for 513 points, \$0.001215/bbl for 1,025
points, and \$0.000275/bbl for 2,049 points. The largest
absolute PI difference between 513 and 4,097 grid points was
approximately $4.5\times10^{-9}$ USD/bbl. Thus coverage
\emph{classifications} are stable in this bounded design, but
individual endpoints are not certified to sub-cent precision.
The observed nominal 95\% coverage under $\rho=1$ is a
finite-sample synthetic frequency, not a validated statement about
actual WTI market-price intervals. This is a pilot-informed
experiment, not a preregistered external confirmation.

\section*{S6. Market-price provenance and interpretation}
The empirical response is the Barchart end-of-day \texttt{Latest}
field. Its equivalence to a final exchange settlement observation,
the first-publication time, and the historical revision vintage
have not been independently verified for the input records.
The LO comparison uses earlier-date option histories, but
repricing uses contemporaneous curve inputs and known fixing
states. In addition, CME Clearing Advisory 20-198 dated
May 15, 2020 states that a Curran model became the end-of-day
settlement valuation model for WTI Average Price Options effective
May 18, 2020. This creates a possible shared-method dependence
between the evaluation mark and the Curran repricing proxy.

\noindent\textbf{Source note.}
All tables preserve results from the manuscript's versioned
research artifacts. The numerical grid sensitivity uses script
\texttt{experiments/jfm\_grid\_sensitivity.py} and the original
independent-seed data, without running new payoff simulations.
Raw licensed contract-level market data are not redistributed.

\end{document}
